\documentclass[11pt]{article}
\usepackage{amsmath}
\usepackage{amssymb}
\usepackage{geometry}
\geometry{margin=1in}

\title{Plan 01-07: Combined Gravitational Potential}
\author{J2+J3+Drag Secular Perturbation Theory}
\date{2026-04-02 — Phase 1}

\begin{document}
\maketitle

\section{Objective}
Combine J2 and J3 disturbing potentials into total gravitational perturbation function.

\section{Combined Gravitational Disturbing Potential}

\begin{equation}
\boxed{R_{\text{grav}}(a,e,i,\Omega,\omega,M) = R_{J2} + R_{J3}}
\label{eq:rgrav}
\end{equation}

Explicitly:
\begin{align}
R_{J2} &= \frac{\mu J_2 R_E^2}{4r^3(M)}\left[(3\sin^2 i - 2) - 3\sin^2 i \cos 2(\omega + f(M))\right]\\
R_{J3} &= \frac{\mu J_3 R_E^3}{r^4(M)} P_3[\sin i \sin(\omega + f(M))]
\end{align}

where:
\begin{itemize}
    \item $r = a(1 - e\cos E(M))$ via Kepler equation $M = E - e\sin E$
    \item $f = f(E)$ via true-to-eccentric anomaly relation
    \item All dependencies on $M$ are explicit through $E(M)$ and $f(E)$
\end{itemize}

\section{Dependency Structure}

\begin{center}
\begin{tabular}{c|l}
\textbf{Element} & \textbf{How it enters $R_{\text{grav}}$} \\
\hline
$a$ & Through $r \propto a$ (both $r^{-3}$ and $r^{-4}$ terms) \\
$e$ & Through $r(e,E)$ and $f(e,E)$ (series expansions) \\
$i$ & Explicit: $\sin^2 i$ (J2), $\sin i$ (J3) \\
$\Omega$ & \textbf{ABSENT} (axisymmetry) \\
$\omega$ & Explicit: $(\omega + f)$ in trigonometric terms \\
$M$ & Through $E(M)$ via Kepler equation \\
\end{tabular}
\end{center}

\section{Axisymmetry Verification}

Both J2 and J3 potentials are axisymmetric (no $\lambda$ or $\Omega$ dependence):

\begin{equation}
\frac{\partial R_{\text{grav}}}{\partial \Omega} = \frac{\partial R_{J2}}{\partial \Omega} + \frac{\partial R_{J3}}{\partial \Omega} = 0 + 0 = 0
\end{equation}

\textbf{Consequence:} From Lagrange equations, this leads to $\langle di/dt \rangle_{\text{secular}} = 0$ (no secular change in inclination from gravitational harmonics alone).

\section{Magnitude Comparison}

At LEO (h = 700 km, r $\approx$ 7078 km):
\begin{align}
\left|\frac{R_{J2}}{U_{\text{point}}}\right| &\approx J_2\left(\frac{R_E}{r}\right)^2 \approx 8.8 \times 10^{-4}\\
\left|\frac{R_{J3}}{U_{\text{point}}}\right| &\approx J_3\left(\frac{R_E}{r}\right)^3 \approx 2.0 \times 10^{-6}
\end{align}

J3 is ~0.23\% of J2 — small but non-negligible for long-term predictions.

\section{Structure Ready for Phase 3}

$R_{\text{grav}}$ is now fully expressed in orbital elements. In Phase 3, we will:
\begin{enumerate}
    \item Compute partial derivatives: $\partial R_{\text{grav}}/\partial a$, $\partial R_{\text{grav}}/\partial e$, etc.
    \item Expand in Fourier series in $M$ using anomaly conversions from Plan 01-06
    \item Time-average: $\langle \partial R_{\text{grav}}/\partial x \rangle$
    \item Substitute into Lagrange equations to get secular rates
\end{enumerate}

\section{Success Criteria}
\begin{itemize}
    \item[$\checkmark$] Combined $R_{\text{grav}} = R_{J2} + R_{J3}$ written
    \item[$\checkmark$] All dependencies documented
    \item[$\checkmark$] Axisymmetry verified ($\partial R/\partial \Omega = 0$)
    \item[$\checkmark$] Magnitude comparison provided
    \item[$\checkmark$] Framework ready for Phase 3
\end{itemize}

\end{document}
